\documentclass[10pt,a4paper]{amsart}
\usepackage[margin=19mm]{geometry}
\usepackage{amsmath,amssymb,mathtools}
\usepackage[scr=rsfs,cal=euler]{mathalfa}
\usepackage{xfrac}
\usepackage[unicode,hidelinks,bookmarks=false]{hyperref}
\newcommand{\ie}{i.e.\ }
\newcommand{\cf}{cf.\ }
\newcommand{\resp}{respectively\ }
\newcommand{\Subsection}{\subsection}
\providecommand{\nobreakdash}{\nobreak\hbox{-}}
\setlength{\emergencystretch}{2em}
\multlinegap0pt
\usepackage{bm}
\usepackage{bbm}
\usepackage{dsfont}
\usepackage{xspace}
\usepackage{cancel}
\usepackage{mathtools}
\usepackage{amsthm}
\makeatletter
\@ifundefined{thm}{\newtheorem{thm}{Theorem}}{}
\@ifundefined{cor}{\newtheorem{cor}[thm]{Corollary}}{}
\@ifundefined{rem}{\newtheorem{rem}[thm]{Remark}}{}
\makeatother
\title[Composition operators on Hardy-Smirnov spaces]{Composition operators on Hardy-Smirnov spaces}
\author{V.V. F\'avaro, P.V. Hai, D.M. Pellegrino, O.R. Severiano}
\date{Source: arXiv:2111.10609v3 (CC BY 4.0)}
\begin{document}
\maketitle

\noindent\emph{Source and licence.} Composition operators on Hardy-Smirnov spaces. Version arXiv:2111.10609v3 (CC BY 4.0), distributed under the Creative Commons Attribution 4.0 International licence, as recorded in the arXiv licence metadata for this exact version and shown on the abstract page https://arxiv.org/abs/2111.10609v3. Full rights notice: licenses/arxiv-2111.10609-CC-BY-4.0.txt.\par\smallskip

\noindent\textit{Editorial scope.} Counted unit: the Thm whose source label is \texttt{038} in the article (source0.tex lines 963--997, source label \texttt{038}), delivered together with the article's own complete original proof of it (source0.tex lines 999--1114, one \texttt{proof} environment of 116 lines). No mathematics is written, repaired or re-derived: statement and proof are the article's own text line for line. Required context retained uncounted: 023 (lines 645--650); 078 (lines 668--678). Every definition, display and label of those lines is retained unchanged. Omitted from this package and not redistributed: 1--644, 651--667, 679--962, 998--998, 1115--2412. Nothing inside a retained excerpt is omitted or altered: every retained body is delivered exactly as the article's lines are written, so no omission receipt of any kind accompanies this package. Citations: the retained excerpts carry no citation command at all, so no bibliography entry of the article is transcribed and no reference is redistributed. Reference adaptation: none was needed and none was made; every reference inside the delivered unit resolves inside it, so no unresolved reference or citation is produced. Printed slips kept literal: no printed word, symbol, hypothesis or number of the counted statement or of its proof was altered. Layout and class: the article's own class is \texttt{amsart} with packages including amscd, amsmath, latexsym, amsthm, amsfonts, amssymb, graphicx, color, geometry, hyperref, inputenc; the wrapper is the lane's pinned \texttt{amsart} editorial wrapper at 10pt on A4, which restates the theorem-like environments the retained text needs and adds the article's own macro definitions that the retained text needs (none), with extra wrapper packages (bm, bbm, dsfont, xspace, cancel, mathtools). The delivered numbering is the unit's own. Assets: no retained span contains a figure environment, a graphics-inclusion command, a PDF-page inclusion, a table or a TikZ picture; no image file is copied into this package, the package declares no asset dependency and no image file is redistributed with it. Licence: version arXiv:2111.10609v3 of this article is distributed under the Creative Commons Attribution 4.0 International licence, as recorded in the arXiv licence metadata for this exact version and shown on the abstract page https://arxiv.org/abs/2111.10609v3. A full rights notice is delivered at licenses/arxiv-2111.10609-CC-BY-4.0.txt. Characters: every retained excerpt is pure ASCII, so the delivered engine, XeLaTeX with its default Latin Modern text font, reports no missing character.
\begin{cor}
\label{023} If $C_{\Phi}$ is a bounded operator on $H^2(\Omega)$ then $C_{\Phi}^{*}k_{u}=k_{\Phi(u)},$ for
each $u\in\Omega.$
%Let $\Phi$ be an analytic self-map of $\Omega$ such that $C_{\Phi
%}$ is bounded on $H^{2}(\Omega)$. 
\end{cor}

\begin{rem}
\label{078} We frequently use the following fact: if $\tau(z)=(az+b)/(cz+d)$
is a Riemann map with $|c|^{2}=|d|^{2}$ then $b\overline{d}\neq a\overline
{c}.$ Indeed, since $\tau$ is a Riemann map, the equality $|c|^{2}=|d|^{2}$
implies that $c\neq0$ and $d\neq0.$ Hence
\begin{align*}
b\overline{d}-a\overline{c}=\frac{b}{d}|d|^{2}-\frac{a}{c}|c|^{2}%
=|d|^{2}\left(  \frac{b}{d}-\frac{a}{c}\right)  =|d|^{2}\left(  \frac
{bc-ad}{dc}\right)  \neq0.
\end{align*}
\end{rem}

\begin{thm}
\label{038} Let $\tau$ be the linear fractional Riemann map $z\in \mathbb{U} \mapsto (az+b)/(cz+d),\lambda \in \mathbb{C}\backslash\{0\},$ $\Phi$ and $\Phi_{\star}$ be analytic self-maps of $\Omega$
such that $C_{\Phi}$ and $C_{\Phi_{\star}}$ are bounded on $H^{2}(\Omega).$ Then
$C_{\Phi}^{\ast}=\lambda C_{\Phi_{\star}}$ if and only if one of the following cases occurs

\begin{enumerate}
\item If $|c|^{2}=|d|^{2}$ then there is a complex number $r$ such that
\begin{align}\label{035}
\Phi(w)=\overline{\lambda}^{-1}w+r\quad\text{and}\quad\Phi_{\star}(w)=\lambda
w+\frac{(\lambda-1)(|a|^{2}-|b|^{2})+\lambda\overline{r}(b\overline
{d}-a\overline{c})}{\overline{b}d-\overline{a}c}.
\end{align}


\item If $|c|^{2}\neq|d|^{2}$ and $(|a|^{2}-|b|^{2})(|c|^{2}-|d|^{2}%
)=|b\overline{d}-a\overline{c}|^{2}$ then there is a complex number $r$ such
that%
\begin{align}\label{036}
\Phi(w)=\lambda w+\frac{(r-1)(b\overline{d}-a\overline{c})}{|c|^{2}-|d|^{2}}%
\quad\text{and}\quad\Phi_{\star}(w)=\lambda \overline{r}w+\frac{(\lambda
\overline{r}-1)(b\overline{d}-a\overline{c})}{|c|^{2}-|d|^{2}}.
\end{align}


\item If $|c|^{2}\neq|d|^{2}$ and $(|a|^{2}-|b|^{2})(|c|^{2}-|d|^{2}%
)\neq|b\overline{d}-a\overline{c}|^{2}$ then $\lambda=1$ and there is
$r\in\mathbb{C}$ such that%
\begin{align}\label{043}
\Phi(w)=rw+\frac{(r-1)(b\overline{d}-a\overline{c})}{|c|^{2}-|d|^{2}}%
\quad\text{and}\quad\Phi_{\star}(w)=\overline{r}w+\frac{(\overline
{r}-1)(b\overline{d}-a\overline{c})}{|c|^{2}-|d|^{2}}.
\end{align}

\end{enumerate}
\end{thm}

\begin{proof}
By Corollary \ref{023}, we have that $C_{\Phi}^{\ast}=\lambda C_{\Phi_{\star}%
}$ for some complex number $\lambda\neq0$ if and only if
\begin{align*}
k_{\Phi(u)}(w)=\lambda k_{u}(\Phi_{\star}(w)),\quad u,w\in\Omega
\end{align*}
or equivalently
%by Lemma \ref{024} that}
\begin{align}\label{031}
&  \lambda\left\{  |a|^{2}-|b|^{2}+(b\overline{d}-a\overline{c})\overline
{\Phi(u)}+[(|c|^{2}-|d|^{2})\overline{\Phi(u)}+\overline{b}d-\overline
{a}c]w\right\} \\
&  =|a|^{2}-|b|^{2}+(b\overline{d}-a\overline{c})\overline{u}+[(|c|^{2}%
-|d|^{2})\overline{u}+\overline{b}d-\overline{a}c]\Phi_{\star}(w),\quad
u,w\in\Omega.\nonumber
\end{align}
Now we study the cases, $|c|^{2}=|d|^{2}$ and $|c|^{2}\neq|d|^{2},$ separately.

\textbf{Case 1:} $|c|^{2}=|d|^{2}.$ \newline In this case, it turns out that
\eqref{031} is
\begin{align}\label{032}
&  \overline{\lambda}[  |a|^{2}-|b|^{2}+(\overline{b}d-\overline
{a}c)\Phi(u)+(b\overline{d}-a\overline{c})\overline{w}] \\
&  =|a|^{2}-|b|^{2}+(\overline{b}d-\overline{a}c)u+(b\overline{d}%
-a\overline{c})\overline{\Phi_{\star}(w)},\quad u,w\in\Omega.\nonumber
\end{align}
Differentiating \eqref{032} with respect to $u,$ we obtain
\begin{align}\label{033}
\overline{\lambda}\Phi^{\prime}(u)(\overline{b}d-\overline{a}c)=(\overline
{b}d-\overline{a}c),\quad u\in\Omega.
\end{align}
By Remark \ref{078}, we have $\overline{b}d\neq\overline{a}c$ and hence
\eqref{033} gives $\Phi(u)=\overline{\lambda}^{-1}u+r$ for some complex number
$r.$ Taking complex conjugate in \eqref{032} and differentiating with respect
to $w,$ we obtain $\Phi_{\star}^{\prime}(w)=\lambda.$ Hence, $\Phi_{\star
}(w)=\lambda w+s$ for some complex number $s.$ Replacing the expressions of
$\Phi(u)$ and $\Phi_{\star}(w)$ in \eqref{032}, we have
\begin{align*}
&  \overline{\lambda}(|a|^{2}-|b|^{2})+(\overline{b}d-\overline{a}%
c)u+(\overline{b}d-\overline{a}c)\overline{\lambda}r+(b\overline{d}%
-a\overline{c})\overline{\lambda w}\\
&  =|a|^{2}-|b|^{2}+(\overline{b}d-\overline{a}c)u+(b\overline{d}%
-a\overline{c})\overline{\lambda w}+(b\overline{d}-a\overline{c})\overline
{s},\quad u,w\in\Omega,
\end{align*}
or equivalently,
\[
(\overline{b}d-\overline{a}c)s
=(\lambda-1)(|a|^{2}-|b|^{2})+\lambda\overline{r}(b\overline
{d}-a\overline{c}).
\]
Thus, we obtain \eqref{035}.

\textbf{Case 2:} $|c|^{2}\neq|d|^{2}.$\newline Differentiating \eqref{031}
with respect to $w,$ we obtain
\[
\lambda\lbrack(|c|^{2}-|d|^{2})\overline{\Phi(u)}+\overline{b}d-\overline
{a}c]=[(|c|^{2}-|d|^{2})\overline{u}+\overline{b}d-\overline{a}c]\Phi_{\star
}^{\prime}(w),\quad u,w\in\Omega.
\]
This implies that
\[
\Phi(u)=ru+\frac{(r-1)(b\overline{d}-a\overline{c})}{|c|^{2}-|d|^{2}},\quad
u\in\Omega,
\]
where $\overline{r}=\Phi_{\star}^{\prime}(w)/\lambda.$ Taking the complex
conjugate in \eqref{031} and differentiating with respect to $u,$ we get that
\[
\Phi_{\star}(w)=\lambda\overline{r}w+\frac{(\lambda\overline{r}-1)(b\overline
{d}-a\overline{c})}{|c|^{2}-|d|^{2}},\quad w\in\Omega.
\]
Conversely, if $|c|^2=|d|^2,$ $\Phi$ and $\Phi_{\star}$ are as in \eqref{035}, then a straightforward computation shows that \eqref{031} holds. If $|c|^2\neq |d|^2,$ $\Phi$ and $\Phi_{\star}$ are as in \eqref{036}, then the left side of \eqref{031} is
%Using the expression of $\Phi(u),$ we obtain
%\begin{align*}
%(|c|^{2}-|d|^{2})\overline{\Phi(u)}+\overline{b}d-\overline{a}c  &
%=(|c|^{2}-|d|^{2})\left[  \overline{ru}+\frac{(\overline{r}-1)(\overline
%{b}d-\overline{a}c)}{|c|^{2}-|d|^{2}}\right]  +\overline{b}d-\overline{a}c\\
%&  =\overline{ru}(|c|^{2}-|d|^{2})+(\overline{r}-1)(\overline{b}d-\overline
%{a}c)+\overline{b}d-\overline{a}c\\
%&  =\overline{ru}(|c|^{2}-|d|^{2})+\overline{r}(\overline{b}d-\overline{a}c)
%\end{align*}
%and
%\[
%(b\overline{d}-a\overline{c})\overline{\Phi(u)}=(b\overline{d}-a\overline
%{c})\left[  \frac{\overline{ru}(|c|^{2}-|d|^{2})+(\overline{r}-1)(\overline
%{b}d-\overline{a}c)}{|c|^{2}-|d|^{2}}\right]  .
%\]
%Hence, the left side of \eqref{031} is
\begin{align}\label{0108}
\lambda\left\{  |a|^{2}-|b|^{2}+(|c|^{2}-|d|^{2})\overline{ru}w+\overline
{r}[(b\overline{d}-a\overline{c})\overline{u}+(\overline{b}d-\overline
{a}c)w]+\frac{(\overline{r}-1)|b\overline{d}-a\overline{c}|^{2}}%
{|c|^{2}-|d|^{2}}\right\}  .
\end{align}
%Similarly, using the expression of $\Phi_{\star}(w),$ we get that
 while the right
side of \eqref{031} is
\begin{align}\label{0109}
|a|^{2}-|b|^{2}+(|c|^{2}-|d|^{2})\lambda\overline{ru}w+\lambda\overline
{r}[(b\overline{d}-a\overline{c})\overline{u}+(\overline{b}d-\overline
{a}c)w]+\frac{(\lambda\overline{r}-1)|\overline{b}d-\overline{a}c|^{2}%
}{|c|^{2}-|d|^{2}}.
\end{align}
Putting these together, it follows that \eqref{0108} is equal to \eqref{0109} if and only if
\[
\lambda\left[  |a|^{2}-|b|^{2}+\frac{(\overline{r}-1)|b\overline{d}%
-a\overline{c}|^{2}}{|c|^{2}-|d|^{2}}\right]  =|a|^{2}-|b|^{2}+\frac
{(\lambda\overline{r}-1)|b\overline{d}-a\overline{c}|^{2}}{|c|^{2}-|d|^{2}},
\]
or equivalently,
\begin{align}\label{044}
(\lambda-1)\left(  |a|^{2}-|b|^{2}-\frac{|b\overline{d}-a\overline{c}|^{2}%
}{|c|^{2}-|d|^{2}}\right)  =0.
\end{align}
Therefore, \eqref{036} or \eqref{043} occurs if and only if \eqref{044} holds.
\end{proof}

\end{document}
